% Part: turing-machines
% Chapter: undecidability
% Section: decision-problem

\documentclass[../../../include/open-logic-section]{subfiles}

\begin{document}

\olfileid[fr]{tur}{und}{dec}
\olsection{Le problème de la décision}

On dit que la logique du premier ordre est \emph{décidable}
s'il existe une méthode effective pour déterminer si un
!!{sentence} donné est valide. Or il n'existe aucune méthode
de ce genre : le problème consistant à décider la validité
des énoncés du premier ordre est insoluble. Nous supposons fixé un codage effectif des énoncés sur un alphabet fini ; les noms indexés des symboles sont écrits comme des chaînes finies.

Pour établir ce résultat négatif, nous montrerons que le
problème de la décision ne peut être résolu par une machine
de Turing. Autrement dit, aucune machine qui reçoit sur son
ruban un !!{sentence} du premier ordre ne s'arrête toujours
avec la sortie~$1$ ou~$0$ selon que cet !!{sentence} est valide ou
non. D'après la thèse de Church--Turing, toute fonction
effectivement calculable est calculable par machine de Turing.
Si la « fonction de validité » était effectivement calculable,
elle le serait donc par une machine de Turing. Puisqu'elle ne
l'est pas, elle n'est pas non plus effectivement calculable.

Notre stratégie consiste à réduire le problème de l'arrêt
au problème de la décision. Voici ce que cela signifie.
Nous avons démontré que la fonction~$h(e,w)$, qui vaut~$1$
si la machine désignée par~$e$ s'arrête sur l'entrée~$w$
et~$0$ sinon, n'est pas calculable par machine de Turing. Ici $w$ désigne un mot d'entrée fini ; son arrêt est déjà indécidable puisque les entrées unaires du problème précédent en sont un cas particulier.
Nous montrerons que, s'il existait une machine de Turing
décidant la validité des énoncés du premier ordre, nous
pourrions construire une machine calculant~$h$. Puisque
c'est impossible, aucune machine ne décide la validité.

La première étape consiste à montrer que, pour toute
machine de Turing~$M$ et toute entrée~$w$, nous pouvons
décrire effectivement !!a{sentence} $!T(M, w)$ qui représente
les instructions de~$M$ et l'entrée~$w$, ainsi que
!!a{sentence}~$!E(M, w)$ qui exprime que « $M$ finit par
s'arrêter », de telle sorte que :
\begin{quote}
  $\Entails !T(M, w) \lif !E(M,w)$ si et seulement si~$M$
  s'arrête sur l'entrée~$w$.
\end{quote}
L'essentiel de la preuve consistera à décrire ces énoncés
$!T(M, w)$ et~$!E(M, w)$, puis à vérifier que
$!T(M, w) \lif !E(M, w)$ est valide si et seulement si~$M$
s'arrête sur l'entrée~$w$.

\end{document}
